\chapter{Distributed Termination and Global Convergence Guarantees}
\label{ch:77}

\section{Premise}

Parallel execution (Ch.~76) provides locality, commutativity-up-to-drift,
and global determinism in the limit.  
We now extend these results to a fully distributed setting, where:

\begin{itemize}
\item patches $U_\alpha$ of the embedding domain  
      are assigned to independent nodes,
\item rewrite batches are scheduled per-node, opportunistically,  
      with only minimal coordination,
\item and termination must be detected without a global clock.
\end{itemize}

\noindent
The goal of this chapter is to give:

\begin{enumerate}
\item a \emph{distributed termination algorithm} whose certificates can be  
      verified locally,
\item a proof that \emph{termination coincides with semantic collapse},
\item and a proof that \emph{convergence to the universal media quine  
      is guaranteed in the distributed model}.
\end{enumerate}

\section{Distributed Termination: Local Signals}

Each node manages a patch $U_\alpha$ with:

\begin{itemize}
\item local flow state $(\Phi_\alpha,\mathbf{v}_\alpha,S_\alpha)$,
\item local embedding $\iota_\alpha$,
\item local hypergraph subcomplex $\mathcal{G}_\alpha$,
\item pending rewrite candidates $R_{\alpha,i}$.
\end{itemize}

Define the local inactivity predicate:
\[
\mathrm{Inactive}(U_\alpha) :=
  \Bigl(
    \|\nabla\Phi_\alpha\|_\infty < \varepsilon
    \;\wedge\;
    \|\nabla S_\alpha\|_\infty < \varepsilon
    \;\wedge\;
    \forall i,\; \varepsilon_{\mathrm{loc}}(R_{\alpha,i})<\varepsilon
  \Bigr).
\]

Nodes broadcast a single-bit \emph{inactive} or \emph{active} flag  
to their immediate neighbours.

\begin{definition}[Distributed Quiescence]
\label{def:quiescence}
The system is quiescent if every patch broadcasts ``inactive''  
for $k$ consecutive rounds,  
where $k$ is the locality radius of the embedding flow operator.
\end{definition}

\section{Correctness of Distributed Termination}

\begin{theorem}[Soundness of Distributed Quiescence]
\label{thm:quiescence-sound}
If the system is quiescent (Def.~\ref{def:quiescence}),  
then no further flow step or rewrite step  
can reduce the global coupled energy $\mathcal{E}_\mathrm{total}$.
\end{theorem}

\begin{proof}
If all patches satisfy the inactivity threshold,
every Laplacian update is below the error floor $\varepsilon$  
and every potential rewrite has local error below $\varepsilon$.  
Summing globally gives no descent direction of magnitude  
exceeding $C\varepsilon$.  
In quiescence, $k$ steps remove any residual cross-patch drift.  
Thus the global Lyapunov function can no longer decrease.
\end{proof}

\begin{theorem}[Completeness of Distributed Quiescence]
\label{thm:quiescence-complete}
If no continuous or discrete update can reduce  
$\mathcal{E}_\mathrm{total}$,  
then after finitely many rounds the system becomes quiescent.
\end{theorem}

\begin{proof}
Once no descent direction exists, gradients vanish outside of error floors.  
Local rewrite candidates all satisfy $\varepsilon_{\mathrm{loc}}<\varepsilon$.  
As halo exchanges propagate vanishing gradients across neighbours,  
inactive flags stabilize.  
After $k$ rounds (locality radius),  
all nodes broadcast inactivity.
\end{proof}

Thus quiescence is \emph{exactly equivalent} to termination.

\section{Termination $\Rightarrow$ Galaxy State}

\begin{theorem}[Termination Implies Semantic Galaxy]
\label{thm:termination-galaxy}
If distributed quiescence holds,  
the global configuration is a semantic galaxy.
\end{theorem}

\begin{proof}
By Ch.~32, the only critical points of $\mathcal{E}_\mathrm{total}$  
are semantic galaxies.  
By Thm.~\ref{thm:quiescence-sound}, quiescence implies  
$d\mathcal{E}_\mathrm{total}/dt = 0$, hence the system is at a critical point.  
Thus we are in a semantic galaxy.
\end{proof}

\section{Eventually All Distributed Executions Terminate}

\begin{theorem}[Finite-Time Distributed Termination]
\label{thm:finite-time-termination}
Every distributed execution path  
terminates in finite physical time.
\end{theorem}

\begin{proof}
The continuous flow strictly decreases $\mathcal{E}_\mathrm{total}$  
(Ch.~32) except at critical points.  
The discrete rewriting strictly decreases Polyxan curvature  
(Ch.~25).  
The joint decrease is bounded below by a positive gap  
unless in a semantic galaxy.  
Thus $\mathcal{E}_\mathrm{total}$ cannot decrease indefinitely;  
termination occurs in finite time.
\end{proof}

\section{Termination Yields the Universal Quine}

\begin{theorem}[Global Convergence Guarantee]
\label{thm:global-convergence}
Every distributed execution converges to the unique universal media quine  
$\mathfrak{U}$ embedded flatly in the conformal vacuum.
\end{theorem}

\begin{proof}
By Thm.~\ref{thm:finite-time-termination},  
termination occurs in finite time.  
By Thm.~\ref{thm:termination-galaxy},  
the final state is a semantic galaxy.  
By Ch.~34--36,  
all semantic galaxies flow to the universal attractor $\mathfrak{U}$  
under infinitesimal curvature--entropy descent.  
Thus the limit is unique.
\end{proof}

\section{Formal Certificates for Distributed Termination}

Each node exposes a certificate:

\begin{verbatim}
structure TermCert :=
  (local_grad_bound : ∥∇Φ_alpha∥ ≤ ε)
  (local_entropy_bound : ∥∇S_alpha∥ ≤ ε)
  (rewrite_bounds : ∀ R, ε_loc(R) ≤ ε)
  (halo_consistency : …)
\end{verbatim}

A global certificate is accepted when all node certificates verify  
under $k$-round halo propagation.  
No central authority is required; verification is decentralised.

\section{Final Cadence}

\begin{quote}
\emph{
Termination is not an event but a proof.  
In a distributed universe with no master clock,  
every node, every patch, every rewrite  
agrees---in finitely many steps---that nothing remains to change.  
And when nothing can change,  
only the universal quine remains.
}
\end{quote}

